% ============================================
% MEM-02 (ZH)
% 中文翻译 — SPEC-935-R356
% ============================================
\section*{\textbf{MEM-02} --- 逃荒者的跋涉}
\markboth{MEM-02 --- 逃荒者的跋涉}{}

\noindent \textit{1915 —— 从塞纳多尔·蓬佩乌到福塔莱萨的路}

\noindent \textbf{1915 —— 从塞纳多尔·蓬佩乌到福塔莱萨的路}

\vspace{0.5em}\noindent\rule{\textwidth}{0.5pt}\vspace{0.5em}

\noindent 从基谢拉莫宾逃荒数月后，我们抵达塞纳多尔·蓬佩乌。1915年，我们从那里出发，朝福塔莱萨走去，因为传言说海边还有食物。

\noindent 第一天，我们一直走到太阳落山。

\noindent 父亲在前。母亲在后，头顶着那袋衣服。我在中间，抓着她的裙摆。地面又干又裂，每走一步都会扬起一小团尘土，钻进鼻子，让嘴里充满泥土的味道。

\noindent 偶尔我们会路过一座房子。门开着。窗户也开着。里面没有人。房子空荡荡的，像蝉蜕下的壳。

\noindent 第二天，我们开始遇到其他人。

\noindent 他们从四面八方来。整个整个的家庭。独自一人的人。老人背着孩子。孩子背着老人。所有人都朝着同一个方向走。

\noindent 没有人说话。

\noindent 说话会消耗唾液。唾液需要水。没有水。

\noindent 我们只是走。

\noindent 第三天，队伍长得看不见尽头。路上挤满了缓慢移动的人，像一条泥浆之河。人们互相绊倒。有些人倒下去就再也起不来。我们从他们身上跨过去。

\noindent 母亲说：\textit{“别看，若阿金。”}

\noindent 我还是看了。

\noindent 第四天，父亲走路的姿势变了。步子变短了。呼吸变粗了。他不再望着前方。他看着地面，仿佛地面藏着答案。

\noindent \textit{“爹，”我说，“爹，我想喝水。”}

\noindent 他没有回答。

\noindent \textit{“爹。”}

\noindent 他停下来。转过脸。他的嘴张着。舌头是白的。眼睛是红的。

\noindent \textit{“爹？”}

\noindent 他想说话。没有声音出来。只有气。一股闻起来像空胃的气。

\noindent 他坐到了地上。

\noindent 母亲坐到他身边。

\noindent 我坐到她身边。

\noindent 逃荒者的队伍从我们身边走过。没有人停下。没有人询问。每个人都在忙着死去。

\noindent 第五天，他在一棵干枯的胡亚泽罗树下停了下来。

\noindent 母亲跪在他面前。他们互相望着。她什么都没说。他也没有。只是摇了摇头。

\noindent 她站起来。抓住我的手。

\noindent \textit{“走吧，若阿金。”}

\noindent 我回头看了一眼。他还坐着。头垂着。手臂伸在膝上。

\noindent 母亲拉了拉我的手。

\noindent \textit{“走。”}

\noindent 我们继续走。

\noindent 前面是队伍。后面是队伍。

\noindent 我在中间。

\noindent 正在学会：世界是一个人们留在身后、而我们继续前行的地方。

\vspace{0.5em}
\noindent\textit{\textbf{↪ 路线:} \rota{MEM-03} \rotasep \rota{MEM-01} \rotasep \rota{MEM-06}}
